\documentclass[aspectratio=169,10pt]{beamer}

\usetheme{Madrid}
\usecolortheme{default}

\usepackage{booktabs}
\usepackage{array}

\definecolor{pgoBlue}{RGB}{37,84,150}
\definecolor{pgoGreen}{RGB}{54,130,91}
\definecolor{pgoOrange}{RGB}{191,104,43}

\setbeamercolor{structure}{fg=pgoBlue}
\setbeamertemplate{navigation symbols}{}
\setbeamertemplate{caption}[numbered]

\newcommand{\case}[1]{\texttt{\detokenize{#1}}}
\newcommand{\pct}[1]{#1\%}
\newcommand{\good}[1]{\textcolor{pgoGreen}{#1}}
\newcommand{\cost}[1]{\textcolor{pgoOrange}{#1}}

\title[Dragon and Bunny Static FEM]{Static FEM Comparison on Conservative Cubic Domains}
\subtitle{Dragon and Bunny: Cubic-Linear vs. Tricubic Hermite}
\author{libpgo experiments}
\date{\today}

\begin{document}

\begin{frame}
  \titlepage
\end{frame}

\begin{frame}{Question}
  \begin{block}{Main question}
    On the same conservative cubic domain, does tricubic Hermite produce a visible-surface static displacement closer to a same-domain tet-linear numerical reference than cubic-linear FEM?
  \end{block}

  \vspace{0.4em}
  \begin{itemize}
    \item Compare original cubic-linear, refined cubic-linear, and cubic-tricubic-Hermite.
    \item Use an x5-scale tet solve as the numerical reference.
    \item Add a 2x2x2 refined cubic-linear case as a stronger low-order baseline.
  \end{itemize}
\end{frame}

\begin{frame}{Asset Pipeline}
  \small
  \begin{enumerate}
    \item Start from the original visible surface mesh: \texttt{dragon.obj} or \texttt{bunny.obj}.
    \item Build an r15 conservative cubic domain with the cubic mesher.
    \item Extract the boundary surface of that cubic domain.
    \item Generate a tet mesh from the extracted cubic-domain boundary.
    \item Split every cubic element into 2x2x2 sub-cubes for \case{cubic_linear_x8}.
    \item Apply the same material, gravity, surface attachment, Newton settings, and surface-error evaluation across cases.
  \end{enumerate}

  \vspace{0.7em}
  \begin{block}{Meaning of ``conservative''}
    The surface is voxelized on a cubic grid. A cube is kept if its center is
    inside the input surface, or if the cube intersects the input surface.
    This keeps the surface embedded inside the volume mesh. Otherwise surface
    embedding can fall outside the simulated domain and require extrapolation.
    The tradeoff is extra volume outside the original surface.
  \end{block}
\end{frame}

\begin{frame}{Same-Domain Design}
  \small
  \begin{itemize}
    \item The tet reference is generated from the cubic domain boundary, not directly from the original surface.
    \item Therefore tet, cubic-linear, refined cubic-linear, and Hermite use the same physical volume domain.
    \item The original surface mesh is still used for attachment and for visible-surface error measurement.
  \end{itemize}

  \vspace{0.6em}
  \begin{table}
    \centering
    \begin{tabular}{lcc}
      \toprule
      model & conservative/original volume & fixed surface vertices \\
      \midrule
      dragon & 1.733817x & 289 \\
      bunny & 1.479166x & 18 \\
      \bottomrule
    \end{tabular}
  \end{table}
\end{frame}

\begin{frame}{Cases}
  \small
  \begin{table}
    \centering
    \begin{tabular}{llp{0.45\linewidth}}
      \toprule
      \normalfont case & formulation & role \\
      \midrule
      \case{tet_ref} & tet-linear & same-domain numerical reference \\
      \case{cubic_linear} & cubic-linear & original r15 low-order cubic baseline \\
      \case{cubic_linear_x8} & cubic-linear & same domain; each cube split 2x2x2 \\
      \case{cubic_hermite} & cubic-tricubic-Hermite & high-order method under test \\
      \bottomrule
    \end{tabular}
  \end{table}

  \vspace{0.5em}
  \begin{block}{Why include \case{cubic_linear_x8}?}
    It checks whether the Hermite gain is only a DOF-count effect. The x8 mesh is not exactly equal-DOF, but it is a comparable-DOF low-order baseline.
  \end{block}
\end{frame}

\begin{frame}{DOF Scale}
  \small
  \begin{table}
    \centering
    \begin{tabular}{lrrrr}
      \toprule
      model & tet ref & cubic-linear & cubic-linear x8 & Hermite \\
      \midrule
      dragon & 624,939 & 15,768 & 105,027 & 126,144 \\
      bunny & 363,759 & 9,102 & 62,355 & 72,816 \\
      \bottomrule
    \end{tabular}
  \end{table}

  \vspace{0.6em}
  \begin{table}
    \centering
    \begin{tabular}{lcc}
      \toprule
      model & tet/Hermite DOF ratio & x8-linear/Hermite DOF ratio \\
      \midrule
      dragon & 4.95x & 0.83x \\
      bunny & 5.00x & 0.86x \\
      \bottomrule
    \end{tabular}
  \end{table}
\end{frame}

\begin{frame}{Solver and Metrics}
  \small
  \begin{columns}[T]
    \begin{column}{0.42\linewidth}
      \textbf{Shared solver setup}
      \begin{itemize}
        \item Stable Neo-Hookean material
        \item Gravity: $(0, -9.81, 0)$
        \item Surface attachment coefficient: $10^5$
        \item Max Newton iterations: 300
        \item Gradient tolerance: $10^{-5}$
      \end{itemize}
    \end{column}
    \begin{column}{0.56\linewidth}
      \textbf{Reported metrics}
      \[
        E_{\mathrm{rel}} =
        \frac{\|u - u_{\mathrm{ref}}\|_{2,F}}
             {\|u_{\mathrm{ref}}\|_{2,F}}
      \]
      \[
        E_y =
        \frac{\|u_y - u_{\mathrm{ref},y}\|_{2,F}}
             {\|u_{\mathrm{ref},y}\|_{2,F}}
      \]
      \[
        E_{95} = \operatorname{P95}_{i \in F}
        \|u_i - u_{\mathrm{ref},i}\|_2
      \]
    \end{column}
  \end{columns}
\end{frame}

\begin{frame}{Scale and Constraint Metrics}
  \small
  Let $S$ be all original surface vertices, $P$ attached vertices, and
  $F = S \setminus P$ free surface vertices.

  \vspace{0.6em}
  \[
    \mathrm{RMS}(u) =
    \sqrt{\frac{1}{|F|}\sum_{i \in F}\|u_i\|_2^2}
  \]

  \[
    \Delta_{\mathrm{RMS}} =
    \frac{\mathrm{RMS}(u) - \mathrm{RMS}(u_{\mathrm{ref}})}
         {\mathrm{RMS}(u_{\mathrm{ref}})}
  \]

  \[
    E_{\mathrm{pin}} = \max_{i \in P}\|u_i\|_2
  \]

  \vspace{0.6em}
  Wall time and peak memory are reported as implementation cost metrics.
\end{frame}

\begin{frame}{Dragon: Accuracy}
  \small
  \begin{table}
    \centering
    \begin{tabular}{lrrrr}
      \toprule
      case & DOFs & free rel L2 & y rel L2 & p95 error \\
      \midrule
      \case{cubic_linear} & 15,768 & \pct{13.73} & \pct{13.25} & 0.02643 \\
      \case{cubic_linear_x8} & 105,027 & \pct{2.72} & \pct{2.70} & 0.00527 \\
      \case{cubic_hermite} & 126,144 & \good{\pct{0.69}} & \good{\pct{0.47}} & \good{0.00137} \\
      \bottomrule
    \end{tabular}
  \end{table}

  \vspace{0.8em}
  \begin{itemize}
    \item Refined cubic-linear removes most of the original low-order error.
    \item Hermite still lowers free-surface relative L2 by about 4.0x versus x8 cubic-linear.
  \end{itemize}
\end{frame}

\begin{frame}{Dragon: Cost}
  \small
  \begin{table}
    \centering
    \begin{tabular}{lrrrr}
      \toprule
      case & iterations & wall time & peak memory & pin residual \\
      \midrule
      \case{tet_ref} & 35 & 409.1 s* & 18.33 GiB* & 0.000529 \\
      \case{cubic_linear} & 34 & 20.8 s & 0.57 GiB & 0.000925 \\
      \case{cubic_linear_x8} & 33 & 1,773.4 s & 2.10 GiB & 0.001002 \\
      \case{cubic_hermite} & 35 & \cost{9,502.8 s} & \cost{7.87 GiB} & 0.000548 \\
      \bottomrule
    \end{tabular}
  \end{table}

  \vspace{0.8em}
  \begin{itemize}
    \item Hermite is more accurate but 5.4x slower than x8 cubic-linear in the local runs.
    \item The cost grows faster than DOF count because the Hermite Hessian has denser coupling.
    \item *The x5 tet reference was run on a server, so its cost is not directly comparable.
  \end{itemize}
\end{frame}

\begin{frame}{Bunny: Accuracy}
  \small
  \begin{table}
    \centering
    \begin{tabular}{lrrrr}
      \toprule
      case & DOFs & free rel L2 & y rel L2 & p95 error \\
      \midrule
      \case{cubic_linear} & 9,102 & \pct{19.61} & \pct{19.78} & 0.00542 \\
      \case{cubic_linear_x8} & 62,355 & \pct{6.97} & \pct{6.22} & 0.00195 \\
      \case{cubic_hermite} & 72,816 & \good{\pct{1.06}} & \good{\pct{0.66}} & \good{0.00031} \\
      \bottomrule
    \end{tabular}
  \end{table}

  \vspace{0.8em}
  \begin{itemize}
    \item Bunny uses the same r15 conservative-domain idea and an ear-tip attachment patch.
    \item Hermite lowers free-surface relative L2 by about 6.6x versus x8 cubic-linear.
  \end{itemize}
\end{frame}

\begin{frame}{Bunny: Cost}
  \small
  \begin{table}
    \centering
    \begin{tabular}{lrrrr}
      \toprule
      case & iterations & wall time & peak memory & pin residual \\
      \midrule
      \case{tet_ref} & 26 & 270.0 s* & 11.62 GiB* & 0.000032 \\
      \case{cubic_linear} & 25 & 20.7 s & 0.32 GiB & 0.000068 \\
      \case{cubic_linear_x8} & 26 & 1,927.0 s & 1.68 GiB & 0.000048 \\
      \case{cubic_hermite} & 30 & \cost{12,268.3 s} & \cost{6.02 GiB} & 0.000034 \\
      \bottomrule
    \end{tabular}
  \end{table}

  \vspace{0.8em}
  \begin{itemize}
    \item Hermite is 6.4x slower and about 3.6x higher peak memory than x8 cubic-linear in the local runs.
    \item The attached-patch residuals are all well below $10^{-3}$.
    \item *The x5 tet reference was run on a server, so its cost is not directly comparable.
  \end{itemize}
\end{frame}

\begin{frame}{Displacement Scale}
  \small
  \begin{table}
    \centering
    \begin{tabular}{lrrrr}
      \toprule
      model & tet ref RMS & cubic-linear $\Delta$ & cubic-linear x8 $\Delta$ & Hermite $\Delta$ \\
      \midrule
      dragon & 0.138372 & -13.32\% & -2.50\% & \good{-0.22\%} \\
      bunny & 0.021106 & -14.24\% & -4.06\% & \good{+0.21\%} \\
      \bottomrule
    \end{tabular}
  \end{table}

  \vspace{0.8em}
  \begin{itemize}
    \item $\Delta$ is relative free-surface RMS difference from the x5 tet reference.
    \item The original cubic-linear runs under-deform on both models.
    \item Hermite is within about 0.25\% of the x5 tet reference in displacement scale.
  \end{itemize}
\end{frame}

\begin{frame}{Cross-Case Summary}
  \small
  \begin{table}
    \centering
    \begin{tabular}{lccc}
      \toprule
      model & cubic-linear & cubic-linear x8 & Hermite \\
      \midrule
      dragon free rel L2 & \pct{13.73} & \pct{2.72} & \good{\pct{0.69}} \\
      bunny free rel L2 & \pct{19.61} & \pct{6.97} & \good{\pct{1.06}} \\
      \bottomrule
    \end{tabular}
  \end{table}

  \vspace{0.8em}
  \begin{table}
    \centering
    \begin{tabular}{lcc}
      \toprule
      model & Hermite error reduction vs x8 & Hermite wall-time multiplier vs x8 \\
      \midrule
      dragon & 4.0x lower error & 5.4x slower \\
      bunny & 6.6x lower error & 6.4x slower \\
      \bottomrule
    \end{tabular}
  \end{table}
\end{frame}

\begin{frame}{Interpretation}
  \small
  \begin{itemize}
    \item The x8 cubic-linear case is important: refinement alone explains much of the gap from the original cubic-linear mesh.
    \item Hermite still wins against the comparable-DOF low-order baseline on both geometries.
    \item The result is consistent across these two static cases.
    \item The practical tradeoff is clear: Hermite buys accuracy with substantially higher factorization time and memory.
  \end{itemize}
\end{frame}

\begin{frame}{Limitations}
  \small
  \begin{itemize}
    \item The tet reference is a high-resolution numerical reference.
    \item Only two static load cases are reported.
    \item The conservative domains are larger than the original surfaces: 1.733817x for dragon and 1.479166x for bunny.
    \item \case{cubic_linear_x8} is comparable-DOF, not exactly equal-DOF, relative to Hermite.
    \item Wall times are local measurements and include implementation-specific sparse solver behavior.
  \end{itemize}
\end{frame}

\begin{frame}{Conclusion}
  \small
  \begin{block}{Main result}
    On both dragon and bunny, tricubic Hermite is closer to the same-domain tet-linear numerical reference than both original cubic-linear and 2x2x2 refined cubic-linear.
  \end{block}

  \vspace{0.7em}
  \begin{itemize}
    \item Dragon: free rel L2 drops from \pct{13.73} to \pct{2.72} to \good{\pct{0.69}}.
    \item Bunny: free rel L2 drops from \pct{19.61} to \pct{6.97} to \good{\pct{1.06}}.
    \item The refined cubic-linear case makes the comparison stricter and the Hermite advantage more meaningful.
  \end{itemize}
\end{frame}

\begin{frame}{Reproducibility Pointers}
  \scriptsize
  \begin{itemize}
    \item Dragon report: \texttt{examples/scripts/tricubic-hermit-static-compare/report/dragon\_static\_compare\_conservative\_r15.md}
    \item Bunny report: \texttt{examples/scripts/tricubic-hermit-static-compare/report/bunny\_static\_compare\_conservative\_r15\_ear\_tip.md}
    \item Static compare runner: \texttt{examples/scripts/tricubic-hermit-static-compare/static\_compare.py}
    \item Cubic asset generator: \texttt{examples/scripts/tricubic-hermit-static-compare/generate\_cubic\_mesh.py}
    \item x8 cubic mesh script: \texttt{examples/scripts/tricubic-hermit-static-compare/subdivide\_cubic\_mesh.py}
    \item Tet reference tuning script: \texttt{examples/scripts/tricubic-hermit-static-compare/tune\_tet\_reference.py}
  \end{itemize}
\end{frame}

\end{document}
